% ==============================================================================
% CHAPTER 9: CONCLUSION
% ==============================================================================
\section{Conclusion}
\label{sec:conclusion}

This paper takes ``multi-agent collaborative automatic question answering
systems'' as its central theme, and conducts a structured evidence-mapping
review of literature published through July 2026 along
three main threads: \emph{why collaborate, how to collaborate,} and
\emph{when to collaborate}.

\textbf{At the architectural level}, this paper proposes and elaborates a
three-class collaborative architecture taxonomy---\textbf{Flat,
Hierarchical, and Hybrid}---and their nine sub-paradigms, and through a
capability-matrix comparison systematically analyses the applicability
boundaries of each architecture class in representative vertical application
domains including healthcare, legal, and software engineering.

\textbf{At the underlying mechanism level}, this paper provides an in-depth
dissection of three core technology pillars of multi-agent collaboration:
communication protocols, memory management, and collaborative reasoning,
revealing the mutual constraints between these mechanisms and architectural
choices.

\textbf{The most important original contribution} of this paper is the
derivation, through a Design Science Research Methodology (DSRM)
\citep{peffers2007dsrm} from the theoretical survey presented in the
preceding chapters, of a \textbf{two-phase architecture-selection decision
framework oriented towards application scenarios}.
The framework is examined through a 10-case architecture-selection stress
test (7 fully consistent / 1 minor deviation / 2 significant discrepancies).
The comparison illustrates its diagnostic utility while also exposing an
ambiguity zone when several task requirements compete, thereby identifying
clear boundaries for subsequent refinement.

At the engineering-resource level, this paper constructs a practical
reference map of representative development platforms and QA or agent
benchmarks, while keeping these technical and background sources distinct
from the 14 core empirical studies.

Multi-agent collaborative QA technology stands at a critical transition
window from laboratory validation to large-scale real-world deployment.
The four challenges synthesised in this paper---architecture adaptation under
competing requirements, limited process-level evaluation, insufficient
oversight in high-stakes settings, and limited task-level composability---
motivate five research directions: process-aware evaluation, multi-objective
orchestration, continual co-evolution, declarative system composition, and
risk-adaptive human oversight. These directions form a connected research
agenda rather than one-to-one responses to the challenges.
We hope this survey provides a solid theoretical foundation and clear
directional guidance for subsequent research, facilitating the broader
real-world deployment of multi-agent collaborative QA technology in the
era of artificial general intelligence.
